\documentclass{article}

\usepackage[ngerman]{babel}
\usepackage[T1]{fontenc}
\usepackage{lmodern}
\usepackage{amsmath}
\usepackage{amssymb}
\usepackage{hyperref}
\usepackage[margin=2cm]{geometry}
\usepackage{listings}
\usepackage{xcolor}



\lstset{
  language=Python,
  basicstyle=\ttfamily\small,
  keywordstyle=\color{blue},
  commentstyle=\color{gray},
  stringstyle=\color{teal},
  numbers=left,
  numberstyle=\tiny\color{gray},
  breaklines=true,
  frame=single
}

% Link zum Repo git clone 
% Link tubloud: 
% passwort tubcloud : 

\title{Aufgabenblatt 2 \\ Gruppe 41}
\author{
    Cedric Bertow
    \and
    Joshua Berweger
    \and
    Clemens Peter Andreas Eberwien
    \and
    Hee Joon Yoon
}
\date{}

\begin{document}

\maketitle

\textbf{Link zum Repository: }

 

\textbf{Link zum Screen Recording: https://tubcloud.tu-berlin.de/s/BDf7oEZWTfQjy9T}



\textbf{passwort tubcloud: ZWmfKcnQfk} 

\newpage

\section{Report}

\subsection{Vergleich verschiedener CSP-Verfahren zum Lösen von Lollipop-Puzzles}

\subsubsection{Lollipop-Puzzles}

In einem Gitter müssen Lollipops platziert werden.
Ein Lollipop ist ein Kreis (der "Kopf") verbunden mit
einer einzigen geraden Linie (der "Stab") in einer benachbarten Zelle, die in den Kreis mündet.
Im Initialzustand sind bereits Köpfe und Stäbe vorgegeben. Diese müssen zu Lollipops vervollständigt
werden. Dabei dürfen sich Kreise innerhalb einer Zeile oder Spalte nicht gegenseitig "sehen". Ausserdem darf 
ein Lollipop an keinen anderen Lollipop angrenzen. 
Für unser Puzzle haben wir drei verschiedene Gittergrößen ausgewählt: 6x6, 10x10 und 15x15. 
Wir haben uns entschieden ein maximales Gitter der Größe 15x15 zunehmen, da alle weiteren Größen zu unübersichtlich werden und
auch zu lange dauern. 
Jedes Gitter hat zusätzlich noch die Schwierigkeiten Normal und Schwer, die sich insofern unterscheiden, wie viele 
unterschiedliche Lösung es für das Puzzle gibt (Bei schwer gibt es meistens nur bis zu einer richtigen Lösung und bei Normal mehrere).


\newpage
\subsection{Softwaredesign}

\subsubsection{Modellierung von Lollipop-Puzzles als CSP in Python}

\paragraph{Theorie}

In einem Lollipop-Puzzle sind die Zellen des Gitters die Variablen des CSP.
Jeder Zelle muss einer von vier Werten(Leer, Vertikale Linie, Horizontale Linie, Kreis) zugewiesen werden.
Jede Zelle hat die gleichen Constraints mit den Zellen in der gleichen Reihe und denen in der gleichen Spalte des Gitters.

\paragraph{Lollipop-CSP in Python}

In einer Klasse \glqq LollipopGitter\grqq{} speichern wir die relevanten Informationen des Puzzles.
Ein Objekt der Klasse LollipopGitter wird mit angegebener Größe erstellt und erzeugt automatisch alle Variablen, Wertebereiche und Constraints.
\begin{lstlisting}[label={lst:LollipopKonstruktor}]
     def __init__(self, size : int):
        self.size = size
        # Liste der Zellen im Gitter
        self.variables = list()
        for i in range(size):
            for j in range(size):
                self.variables.append((i, j))
        # Dict Mapping aller Zellen auf ihren Wertebereich
        self.domains = {V : list(Domain) for V in self.variables}
        self.constraints = {}
        for variable in self.variables:
            self.constraints[variable] = []
            if variable not in self.domains:
                raise LookupError("Every variable should have a domain assigned to it.")
\end{lstlisting}
Die Variablen werden im gesamten Programm als Tupel aus Integers dargestellt, wobei der erste Wert
die Zeile und der zweite Wert die Spalte der Zelle im Gitter ist.
Der Wertebereich ist für jede Variable eine Liste des Aufzählungstyps Domain.
\begin{lstlisting}[label={lst:domain}]
    class Domain(Enum):
    CIRCLE = 0
    HORIZONTAL_LINE = 1
    VERTICAL_LINE = 2
    EMPTY = 3
\end{lstlisting}

Constraints werden auch als Klasse modelliert.
Für Lollipop Rätsel haben wir zwei verschiedene Arten von Constraints modelliert.
Den AdjacentCellConstraint für Zellen, die orthogonal nebeneinander liegen und den RowColumnConstraint für alle Zellen in
der Reihe und Spalte.
Diese Klassen erben beide von der Klasse Constraint und unterscheiden sich nur in ihren Abfragen.
\begin{lstlisting}[label={lst:ConstraintKlasse}]
    class Constraint(ABC):
    def __init__(self, variables: List):
        self.variables = variables

    @abstractmethod
    def satisfied(self, assignment: Dict) -> bool:
        pass
\end{lstlisting}
Um zu prüfen, ob ein Constraint erfüllt ist, muss eine Zuweisung(assignment) übergeben werden.
Eine Zuweisung ist ein Python Dictionary, welches Variablen auf einen Wert des Wertebereichs abbildet.
Diese Zuweisung wird in jeder Suchfunktion Lokal gespeichert und verändert und nicht im LollipopGitter gespeichert.\newline
LollipopGitter hat außerdem noch zwei wichtige Funktionen.
Die Funktion build\underbar puzzle() wird zum Erstellen eines Rätsels, also einer Startzuweisung benutzt.
Die Funktion consistent() überprüft, ob nach Zuweisung einer neuen Variable, die Zuweisung aller Variablen noch gültig ist.
\begin{lstlisting}[label={lst:consistent}]
  def consistent(self, variable: Tuple[int,int], assignment: Dict) -> bool:
        for constraint in self.constraints[variable]:
            if not constraint.satisfied(assignment):
                return False
        return True
\end{lstlisting}

\subsubsection{Suchfunktionen}
Wir haben folgende vier Suchfunktionen implementiert:
\begin{description}
    \item [Backtracking] Von oben links im Gitter anfangend wird über alle Variablen eine einfache Backtracking Suche ausgeführt.
    \item[Forward-Checking] Struktur analog zu Backtracking, nachdem eine neue Variable belegt wird, werden alle Wertebereiche der umliegenden Zellen gekürzt.
    \item[Forward-Checking mit Heuristik] Unterscheidet sich zum Vorgänger nur in der Reihenfolge der Zellen, welche untersucht werden.
    Nach Minimum-Remaining-Values Heuristik wird immer die Zelle zuerst untersucht, welche am wenigsten mögliche Werte hat.
    \item[Complete Search] Analog zum Vorgänger, nur am Anfang wird über das ganze Gitter mit Constraint Propagation 2-Konsistenz hergestellt.
\end{description}


\subsection{Darstellung der Leistungsmerkmale}
Um die Suchfunktionen zu vergleichen haben wir drei Merkmale:
\begin{description}
    \item[Backtrackingschritte]
    \item[Variablenbelegungen]
    \item[Laufzeit]
\end{description}
Die Complete Search ist in allen Punkten dominant über jedes Puzzle.
Genauso ist die Forward-Checking Suche mit Heuristik den beiden Vorgängern überlegen.
Interessant ist, dass die Forward-Checking-Suche ohne Heuristik zwar weniger Variablenbelegungen prüft als die Backtracking-Suche aber trotzdem höhere Laufzeit hat.
Um die Complete-Search zu optimieren haben wir probiert, Constraint-Propagation in jedem Schritt anzuwenden.
Das Ergebnis war zwar reduzierte Anzahl an Backtrackingschritten, jedoch war die Laufzeit deutlich höher.
Somit ist eine Forward-Checking-Suche mit Heuristik und Constraint-Propagation am Anfang die beste Methode um Lollipop-Puzzles zu lösen.

\subsection{Reflexion}

Den eigentlichen CSP-Solver (csp\_builder.py, csp\_solver.py, puzzel\_lib.py) haben wir bewusst 
ohne KI geschrieben. Nur die Datei visualizer.py haben wir mit Claude Opus auf Medium
erstellen lassen. Sie zeigt die Startkonfiguration und den gesamten Suchpfad inklusive Backtracking 
in einem eigenen Fenster und hat uns so beim Debuggen geholfen. Dafür haben wir Claude unseren 
gesamten Code gegeben und den Visualizer nach jeder Änderung am Solver neu erzeugen lassen. Die KI 
war sehr hilfreich und einfach zu benutzen: Meist entstand in unter zehn Minuten funktionierender Code, 
den wir ohne Debugging einbinden konnten.

\newpage
\section{Reflexion}

Unser Lollipop-Puzzle-Solver besteht aus vier Python-Dateien, die nach und nach entstanden sind: csp\_builder.py, csp\_solver.py, puzzel\_lib.py und visualizer.py. 
Alle außer visualizer.py haben wir selbst geschrieben.
Uns war es wieder sehr wichtig, das Problem ohne KI zu lösen. Um das Projekt ein bisschen schöner zu gestalten und beim Debuggen Hilfe zu kriegen haben wir uns 
dieses hilfreiche Tool erstellen lassen. Es war ganz einfach. Immer wenn wir unseren Code verändert haben, haben wir Claude Opus auf Medium die Aufgabe gegeben 
den Visualizer neu zu erstellen, 
mithilfe unseres gesamten Codes. Es hat meist auf anhieb geklappt. Dafür haben wir folgenden Prompt verwendet (Originalwortlaut):

All unser Code und dann dieser Prompt:

\"
ich will null code von mir besser haben oder verändert
alles was am ende da sein soll ist eine neue python datei namens visualizer
im gleichem Folder namens sec sind auch die anderen beiden dateien
eins namens csp\_builder.py das andere namens csp\_solver.py
dieses soll eine mapp von einem Gitter zeigen mit der start Konfiguration und danach den gesamten Denkpfad also das backtracking am besten mit schönem interface irgendwie was sinn macht
das ganze soll auf einer extra canvis passieren
\"

Dazu sei gesagt, dass wir im Laufe der Zeit mehrere Versionen des Visualizers benötigt haben, die sich aber alle sehr ähnlich waren.

Insgesamt war die KI sehr hilfreich und vor allem einfach zu benutzen. Meist hat sie uns in unter zehn Minuten 
wirklich guten Code erstellt, den wir sofort und ohne Debugging einbinden konnten. 
Den generierten Code haben wir allerdings nicht im Detail nachvollzogen, sondern nur 
geprüft, ob die Visualisierung korrekt ist. Gerade deshalb war es uns wichtig, den eigentlichen Solver selbst zu schreiben.
\end{document}
